% !TeX root = ../../Book3.tex
\section{极值方法}
\label{b3:sec:0502}

存在性的证明不直接解一个映射公式，而是把原点附近的伸缩率推到最大。正规族保证极值能达到；若极值映射漏掉一点，就用平方根重新铺开像，使伸缩率变大。两项共同给出满射性。

\subsection{归一化族}
\label{b3:subsec:050201}

固定真单连通区域 \(U\subsetneq\mathbb C\) 与 \(a\in U\)。令 \(\mathcal F\) 为所有单叶全纯映射 \(f:U\rightarrow\mathbb D\)，满足 \(f(a)=0\)、\(f'(a)>0\)。
\begin{lemma}\label{b3:lem:05-nonempty-normalized}
归一化族 \(\mathcal F\) 非空。
\end{lemma}
\begin{proof}
选 \(b\notin U\)，取 \(h(z)^2=z-b\) 的全纯根。它单叶，像与负像不交。由于像在 \(h(a)\) 附近包含一个小圆盘 \(D(h(a),r)\)，像便避开 \(D(-h(a),r)\)。因此 \(q(z)=1/(h(z)+h(a))\) 全纯单叶且模不超过 \(1/r\)。将 \(q\) 乘一个足够小的正数，使像在开单位盘内；接着用圆盘自同构把 \(a\) 的像送到0，再用旋转使导数正实。所得映射属于 \(\mathcal F\)。
\end{proof}
不能直接取 \(1/(z-b)\) 并声称它有界，因为 \(b\) 可能是边界点，\(z\) 可任意接近它。平方根像与负像的分离提供了真正的正距离，而不只是分母处处非零。

在构造有界单叶函数时，取倒数以前必须留下一个真正被像集避开的圆盘。仅知道某点不在像中，不能保证倒数有界；例如 \(1/z\) 在去心平面无零，倒数仍可无界。平方根的像与负像都是开集，又互不相交。在负像中选定一点及其小邻域，便取得整个小圆盘与正像不交的定量间隔。之后才有 \(|1/(h-c)|\le1/\varepsilon\)。这一开集间隔是 Riemann 构造中最容易被“选一个外点”一句话遮住的地方。

\ProseParagraphBreak
还需区分“函数有界”与“像已经是圆盘”。第一次倒数构造只把区域单射地映入一个有界集合。仿射缩放使它成为圆盘内的单叶函数，最后的极值步骤才把像中的遗漏点消去。这样证明并没有事先假设存在满射圆盘坐标，也没有用边界光滑性制造一个预先的映射。

\subsection{正规族紧性}
\label{b3:subsec:050202}

\begin{lemma}\label{b3:lem:05-extremal-compactness}
数 \(S=\sup_{f\in\mathcal F}f'(a)\) 满足 \(0<S<\infty\)。若 \(f_n\in\mathcal F\) 且 \(f_n'(a)\rightarrow S\)，则有局部一致子列，其极限仍属于 \(\mathcal F\)。
\end{lemma}
\begin{proof}
非空性给 \(S>0\)。选 \(\overline{D(a,r)}\subset U\)，全部成员模小于1，所以 Cauchy 估计给 \(f'(a)\le1/r\)，故 \(S\) 有限。Montel 定理提供局部一致子列极限 \(f\)，且 \(|f|\le1\)、\(f(a)=0\)。导数收敛给 \(f'(a)=S>0\)，故它非常数。最大模原理使 \(|f|<1\)，单叶极限定理使 \(f\) 单叶，导数为正实数，因此 \(f\in\mathcal F\)。
\end{proof}
像的边界条件在极限中最容易丢失。单纯的 \(|f_n|<1\) 只给 \(|f|\le1\)；先排除常数退化，才能用最大模把闭盘目标重新缩成开盘。

极值列的导数并非自动有极限。先由 Cauchy 估计给出其有限上界，选取导数趋于上确界的序列，再用 Montel 提取紧集一致收敛子列；导数收敛来自第4章。归一化点处的极限导数为正，使极限不可能常值。单叶函数的紧极限只有“仍单叶”与“常值”两个选择，故正导数恰好排除了丢失单射性的退化。若只要求 \(f(a)=0\) 而不控制导数，\(f_n=f/n\) 就是趋于常值0的合法单叶列。

\ProseParagraphBreak
把函数类中的导数要求为正实数也不是额外限制。任意非零导数都能靠乘一个单位复数转正，而导数的模保持不变。因此最大化正导数与最大化导数模是同一极值问题。这个归一化的好处是極限导数无需在复平面中追踪变化的方向，唯一性证明也不再留有一个未指定旋转。

极值论证对函数族作紧化，而没有对区域边界作紧化。局部有界性使函数在紧集上不会无界，Cauchy 估计进一步控制它们的振荡；对角抽取才把每个局部窗口的极限拼成整域函数。若删去局部有界性，序列 \(nz\) 无法取得有限全纯极限；若保留有界性却删去正导数归一化，极限可能塌成常值。证明的各条件分别阻止不同的退化，因此不能把 Montel 一词视为一个自动生成所需极值函数的黑盒。

\ProseParagraphBreak
导数的上确界也不是由每个函数在区域全体的最大值给出，因为函数模的上确界可为1且不在内部取到。这里最大化的是一个内点的导数，Cauchy 公式把它表示为固定小圆周上的连续线性泛函，紧集一致收敛便保持其值。这个局部量容易在极限中保存，最终却控制了像是否满圆盘，是极值证明尤其有效的原因。

\subsection{极值函数}
\label{b3:subsec:050203}

前一引理给出达到 \(S\) 的函数 \(f\)。极值对象不是区域的面积、边界长度或全域最大导数，而是固定点 \(a\) 处归一化导数的最大值。固定点与相位归一化使比较数是正实数，避免在复数上说“最大导数”。

\ProseParagraphBreak

极值问题的可达性依赖三个封闭步骤：归一化条件由点值与导数收敛保持，单射条件由 Hurwitz 保持，开圆盘目标由非退化与最大模保持。缺少其中任何一项，正规族只给某个极限，并不保证它仍是候选映射。这个论证也解释为什么选择原点导数作目标，而不是一个可能在边界逃逸的全域上确界。

遗漏点的平方根步骤需要谨慎计算它对归一化导数的作用。设当前圆盘像遗漏 \(b\ne0\)，先令 \(g=\phi_b\circ f\)，然后在源区域上取根 \(h^2=g\)。这时 \(|h|<1\)、\(|h(a)|=\sqrt{|b|}\)，但 \(h(a)\) 不再为0。将它送回0所需的自同构，导数模为 \(1/(1-|h(a)|^2)\)。另一方面，平方根链式法则给 \(|h'(a)|=|g'(a)|/(2|h(a)|)\)。合并而非遗漏其中任一因子，才得到改善比
\[
 \frac{1+|b|}{2\sqrt{|b|}}>1.
\]
当遗漏点很靠近0时，这个比值很大；靠近圆周时改善趋近于1，但仍严格大于1。证明所需的只是严格性，不需统一改善量。

\ProseParagraphBreak
这里根函数保持单射，因为根的相同值平方后给 \(g\) 的相同值，\(g\) 又单射。平方根的分支选哪一个不影响导数模，最后的相位归一化还会消除这两个选择。极值论证因此既不要求遗漏集合有面积，也不要求能在遗漏集合中找到一个圆盘；一个遗漏点已足以产生矛盾。和最开始构造有界函数时需要避开的圆盘相比，两步使用外点的方式完全不同。

\subsection{满射性论证}
\label{b3:subsec:050204}

\begin{lemma}[极值映射的满射性]\label{b3:lem:05-extremal-surjective}
达到最大归一化导数的 \(f\in\mathcal F\) 满足 \(f(U)=\mathbb D\)。
\end{lemma}
\begin{proof}
若漏掉 \(b\in\mathbb D\)，因 \(f(a)=0\) 得 \(b\ne0\)。令 \(F=\phi_b\circ f\)，则 \(F\) 无零。区域单连通，存在全纯根 \(h^2=F\)，且 \(|h|<1\)。它单叶，因为 \(h(z_1)=h(z_2)\) 推出 \(F(z_1)=F(z_2)\)，再推出两点相同。令 \(c=h(a)\)，则 \(|c|=\sqrt{|b|}\)。函数 \(g=\phi_c\circ h\) 单叶、取值于圆盘且 \(g(a)=0\)。其导数的模为
\[
 |g'(a)|=\frac{|h'(a)|}{1-|c|^2}
 =\frac{1-|b|^2}{2\sqrt{|b|}(1-|b|)}\cdot f'(a)
 =\frac{1+|b|}{2\sqrt{|b|}}\cdot f'(a).
\]
中间等式使用 \(2h(a)h'(a)=F'(a)\) 与 \(|\phi_b'(0)|=1-|b|^2\)。由于 \(0<|b|<1\)，有 \(1+|b|>2\sqrt{|b|}\)。旋转 \(g\) 使导数正实，就得到 \(\mathcal F\) 中导数大于 \(S\) 的成员，矛盾。
\end{proof}
平方根把遗漏值变成一个可利用的放大因子。它并不将原映射像直接补齐，而是构造一个新的候选映射；极值矛盾保证原映射其实无遗漏。若 \(|b|=1\)，因子等于1，不会产生矛盾，但边界点本来就不是满射到开盘所要求的点。
